% Options for packages loaded elsewhere
\PassOptionsToPackage{unicode}{hyperref}
\PassOptionsToPackage{hyphens}{url}
\documentclass[
]{article}
\usepackage{xcolor}
\usepackage[margin=1in]{geometry}
\usepackage{amsmath,amssymb}
\setcounter{secnumdepth}{-\maxdimen} % remove section numbering
\usepackage{iftex}
\ifPDFTeX
  \usepackage[T1]{fontenc}
  \usepackage[utf8]{inputenc}
  \usepackage{textcomp} % provide euro and other symbols
\else % if luatex or xetex
  \usepackage{unicode-math} % this also loads fontspec
  \defaultfontfeatures{Scale=MatchLowercase}
  \defaultfontfeatures[\rmfamily]{Ligatures=TeX,Scale=1}
\fi
\usepackage{lmodern}
\ifPDFTeX\else
  % xetex/luatex font selection
\fi
% Use upquote if available, for straight quotes in verbatim environments
\IfFileExists{upquote.sty}{\usepackage{upquote}}{}
\IfFileExists{microtype.sty}{% use microtype if available
  \usepackage[]{microtype}
  \UseMicrotypeSet[protrusion]{basicmath} % disable protrusion for tt fonts
}{}
\makeatletter
\@ifundefined{KOMAClassName}{% if non-KOMA class
  \IfFileExists{parskip.sty}{%
    \usepackage{parskip}
  }{% else
    \setlength{\parindent}{0pt}
    \setlength{\parskip}{6pt plus 2pt minus 1pt}}
}{% if KOMA class
  \KOMAoptions{parskip=half}}
\makeatother
\usepackage{color}
\usepackage{fancyvrb}
\newcommand{\VerbBar}{|}
\newcommand{\VERB}{\Verb[commandchars=\\\{\}]}
\DefineVerbatimEnvironment{Highlighting}{Verbatim}{commandchars=\\\{\}}
% Add ',fontsize=\small' for more characters per line
\usepackage{framed}
\definecolor{shadecolor}{RGB}{248,248,248}
\newenvironment{Shaded}{\begin{snugshade}}{\end{snugshade}}
\newcommand{\AlertTok}[1]{\textcolor[rgb]{0.94,0.16,0.16}{#1}}
\newcommand{\AnnotationTok}[1]{\textcolor[rgb]{0.56,0.35,0.01}{\textbf{\textit{#1}}}}
\newcommand{\AttributeTok}[1]{\textcolor[rgb]{0.13,0.29,0.53}{#1}}
\newcommand{\BaseNTok}[1]{\textcolor[rgb]{0.00,0.00,0.81}{#1}}
\newcommand{\BuiltInTok}[1]{#1}
\newcommand{\CharTok}[1]{\textcolor[rgb]{0.31,0.60,0.02}{#1}}
\newcommand{\CommentTok}[1]{\textcolor[rgb]{0.56,0.35,0.01}{\textit{#1}}}
\newcommand{\CommentVarTok}[1]{\textcolor[rgb]{0.56,0.35,0.01}{\textbf{\textit{#1}}}}
\newcommand{\ConstantTok}[1]{\textcolor[rgb]{0.56,0.35,0.01}{#1}}
\newcommand{\ControlFlowTok}[1]{\textcolor[rgb]{0.13,0.29,0.53}{\textbf{#1}}}
\newcommand{\DataTypeTok}[1]{\textcolor[rgb]{0.13,0.29,0.53}{#1}}
\newcommand{\DecValTok}[1]{\textcolor[rgb]{0.00,0.00,0.81}{#1}}
\newcommand{\DocumentationTok}[1]{\textcolor[rgb]{0.56,0.35,0.01}{\textbf{\textit{#1}}}}
\newcommand{\ErrorTok}[1]{\textcolor[rgb]{0.64,0.00,0.00}{\textbf{#1}}}
\newcommand{\ExtensionTok}[1]{#1}
\newcommand{\FloatTok}[1]{\textcolor[rgb]{0.00,0.00,0.81}{#1}}
\newcommand{\FunctionTok}[1]{\textcolor[rgb]{0.13,0.29,0.53}{\textbf{#1}}}
\newcommand{\ImportTok}[1]{#1}
\newcommand{\InformationTok}[1]{\textcolor[rgb]{0.56,0.35,0.01}{\textbf{\textit{#1}}}}
\newcommand{\KeywordTok}[1]{\textcolor[rgb]{0.13,0.29,0.53}{\textbf{#1}}}
\newcommand{\NormalTok}[1]{#1}
\newcommand{\OperatorTok}[1]{\textcolor[rgb]{0.81,0.36,0.00}{\textbf{#1}}}
\newcommand{\OtherTok}[1]{\textcolor[rgb]{0.56,0.35,0.01}{#1}}
\newcommand{\PreprocessorTok}[1]{\textcolor[rgb]{0.56,0.35,0.01}{\textit{#1}}}
\newcommand{\RegionMarkerTok}[1]{#1}
\newcommand{\SpecialCharTok}[1]{\textcolor[rgb]{0.81,0.36,0.00}{\textbf{#1}}}
\newcommand{\SpecialStringTok}[1]{\textcolor[rgb]{0.31,0.60,0.02}{#1}}
\newcommand{\StringTok}[1]{\textcolor[rgb]{0.31,0.60,0.02}{#1}}
\newcommand{\VariableTok}[1]{\textcolor[rgb]{0.00,0.00,0.00}{#1}}
\newcommand{\VerbatimStringTok}[1]{\textcolor[rgb]{0.31,0.60,0.02}{#1}}
\newcommand{\WarningTok}[1]{\textcolor[rgb]{0.56,0.35,0.01}{\textbf{\textit{#1}}}}
\usepackage{graphicx}
\makeatletter
\newsavebox\pandoc@box
\newcommand*\pandocbounded[1]{% scales image to fit in text height/width
  \sbox\pandoc@box{#1}%
  \Gscale@div\@tempa{\textheight}{\dimexpr\ht\pandoc@box+\dp\pandoc@box\relax}%
  \Gscale@div\@tempb{\linewidth}{\wd\pandoc@box}%
  \ifdim\@tempb\p@<\@tempa\p@\let\@tempa\@tempb\fi% select the smaller of both
  \ifdim\@tempa\p@<\p@\scalebox{\@tempa}{\usebox\pandoc@box}%
  \else\usebox{\pandoc@box}%
  \fi%
}
% Set default figure placement to htbp
\def\fps@figure{htbp}
\makeatother
\setlength{\emergencystretch}{3em} % prevent overfull lines
\providecommand{\tightlist}{%
  \setlength{\itemsep}{0pt}\setlength{\parskip}{0pt}}
\usepackage{bookmark}
\IfFileExists{xurl.sty}{\usepackage{xurl}}{} % add URL line breaks if available
\urlstyle{same}
\hypersetup{
  pdftitle={1.2 Final Report},
  pdfauthor={Thoa Le (a1847933)},
  hidelinks,
  pdfcreator={LaTeX via pandoc}}

\title{1.2 Final Report}
\author{Thoa Le (a1847933)}
\date{}

\begin{document}
\maketitle

Data from 1960 to 2024 on Total Live Births (TLB) and Total Fertility
Rates (TFR) was pulled from Singapore Department of Statistics
(data.gov.sg). The data was then transformed into a tsibble object with
three columns: Year, Total Fertility Rates (TFR) and Total Live Births.

\begin{verbatim}
## [1] TRUE
\end{verbatim}

From the exploratory data analysis, the data was visualised to identify
patterns and trends over time. Since the data is yearly, classical
decomposition and the X-11 method are not applicable and seasonal
patterns cannot be detected. Additionally, the socio-economic and
political factors influencing fertility are not included. Once the
patterns are identified, appropriate models can be selected, fitted, and
evaluated for forecasting.

\newpage

\begin{center}\includegraphics[width=0.7\linewidth]{Report_files/figure-latex/page-2-1} \end{center}

\begin{center}\includegraphics[width=0.7\linewidth]{Report_files/figure-latex/page-2-2} \end{center}

From the initial plots of TFR and TLB, both series show a long-term
decreasing trend over time, with TFR exhibiting a more pronounced
decline. A notable change is observed around 1980 to 1990, where both
series display a temporary broad peak before continuing their downward
trend. As the data is yearly, seasonality cannot be determined. While
both series may suggest cyclical behaviour, the trend would need to be
removed to confirm this.

\newpage

\pandocbounded{\includegraphics[keepaspectratio]{Report_files/Report_files/figure-latex/page-3-1.pdf}}

\begin{verbatim}
## NULL
\end{verbatim}

\pandocbounded{\includegraphics[keepaspectratio]{Report_files/Report_files/figure-latex/page-3-2.pdf}}

\begin{verbatim}
## NULL
\end{verbatim}

By using differencing, the plots look like the trend has been removed.

\newpage

\begin{center}\includegraphics[width=0.7\linewidth]{Report_files/figure-latex/page-4-1} \end{center}

\begin{center}\includegraphics[width=0.7\linewidth]{Report_files/figure-latex/page-4-2} \end{center}

The ACF for TFR shows significant positive autocorrelations from lags 1
to 9, while TLB shows significant positive autocorrelations from lags 1
to 6, with a slow decrease in the ACF as the lag increases. This
suggests that neither series is white noise, as there are multiple
spikes outside the significance bounds.

The presence of a strong trend and slow decrease confirms that both TFR
and TLB are non-stationary series. The absence of a sharp cut-off
suggests that differencing or an autoregressive (AR) model may be
suitable for this series.

\begin{Shaded}
\begin{Highlighting}[]
\NormalTok{BirthAndFertility }\SpecialCharTok{|\textgreater{}}
  \FunctionTok{ACF}\NormalTok{(}\StringTok{\textasciigrave{}}\AttributeTok{Total Fertility Rate (TFR)}\StringTok{\textasciigrave{}}\NormalTok{, }\AttributeTok{lag\_max =} \DecValTok{24}\NormalTok{) }\SpecialCharTok{|\textgreater{}}
  \FunctionTok{autoplot}\NormalTok{() }\SpecialCharTok{+}
  \FunctionTok{labs}\NormalTok{(}\AttributeTok{title =} \StringTok{"ACF Plot for Total Fertility Rate (TFR)"}\NormalTok{)}
\end{Highlighting}
\end{Shaded}

\pandocbounded{\includegraphics[keepaspectratio]{Report_files/Report_files/figure-latex/page-10-1.pdf}}
An extended ACF plot for TFR confirms that significant positive
autocorrelations persist up to lag 10.

\newpage

\begin{Shaded}
\begin{Highlighting}[]
\NormalTok{dcmptfr }\OtherTok{\textless{}{-}}\NormalTok{ BirthAndFertility }\SpecialCharTok{|\textgreater{}}
  \FunctionTok{model}\NormalTok{(}\AttributeTok{stl =} \FunctionTok{STL}\NormalTok{(}\StringTok{\textasciigrave{}}\AttributeTok{Total Fertility Rate (TFR)}\StringTok{\textasciigrave{}}\NormalTok{))}
\FunctionTok{components}\NormalTok{(dcmptfr) }\SpecialCharTok{|\textgreater{}}
  \FunctionTok{as\_tsibble}\NormalTok{() }\SpecialCharTok{|\textgreater{}}
  \FunctionTok{autoplot}\NormalTok{(}\StringTok{\textasciigrave{}}\AttributeTok{Total Fertility Rate (TFR)}\StringTok{\textasciigrave{}}\NormalTok{, }\AttributeTok{colour =} \StringTok{"grey"}\NormalTok{) }\SpecialCharTok{+} \CommentTok{\# raw data}
  \FunctionTok{geom\_line}\NormalTok{(}\FunctionTok{aes}\NormalTok{(}\AttributeTok{y =}\NormalTok{ trend), }\AttributeTok{colour =} \StringTok{"\#D55E00"}\NormalTok{) }\SpecialCharTok{+} \CommentTok{\# trend{-}cycle component}
  \FunctionTok{labs}\NormalTok{(}\AttributeTok{title =} \StringTok{"Plot of trend{-}cycle component (orange) and raw data (grey) for TFR"}\NormalTok{)}
\end{Highlighting}
\end{Shaded}

\pandocbounded{\includegraphics[keepaspectratio]{Report_files/Report_files/figure-latex/page-13-1.pdf}}

\newpage

\begin{Shaded}
\begin{Highlighting}[]
\NormalTok{dcmptlb }\OtherTok{\textless{}{-}}\NormalTok{ BirthAndFertility }\SpecialCharTok{|\textgreater{}}
  \FunctionTok{model}\NormalTok{(}\AttributeTok{stl =} \FunctionTok{STL}\NormalTok{(}\StringTok{\textasciigrave{}}\AttributeTok{Total Live{-}Births}\StringTok{\textasciigrave{}}\NormalTok{))}
\FunctionTok{components}\NormalTok{(dcmptlb) }\SpecialCharTok{|\textgreater{}}
  \FunctionTok{as\_tsibble}\NormalTok{() }\SpecialCharTok{|\textgreater{}}
  \FunctionTok{autoplot}\NormalTok{(}\StringTok{\textasciigrave{}}\AttributeTok{Total Live{-}Births}\StringTok{\textasciigrave{}}\NormalTok{, }\AttributeTok{colour =} \StringTok{"grey"}\NormalTok{) }\SpecialCharTok{+} \CommentTok{\# raw data}
  \FunctionTok{geom\_line}\NormalTok{(}\FunctionTok{aes}\NormalTok{(}\AttributeTok{y =}\NormalTok{ trend), }\AttributeTok{colour =} \StringTok{"\#D55E00"}\NormalTok{) }\SpecialCharTok{+} \CommentTok{\# trend{-}cycle component}
  \FunctionTok{labs}\NormalTok{(}\AttributeTok{title =} \StringTok{"Plot of trend{-}cycle component (orange) and raw data (grey) for TLB"}\NormalTok{)}
\end{Highlighting}
\end{Shaded}

\pandocbounded{\includegraphics[keepaspectratio]{Report_files/Report_files/figure-latex/page-14-1.pdf}}
The STL decomposition suggests that both the TFR and TLB follow an
additive structure with a clear trend component evident with the orange
lines. The trend-component follows the TFR raw data smoothly, with
little variability whereas the trend-component of TLB shows greater
variability and larger fluctuations compared to its raw data.

\newpage

\begin{Shaded}
\begin{Highlighting}[]
\FunctionTok{components}\NormalTok{(dcmptfr) }\SpecialCharTok{|\textgreater{}} \FunctionTok{autoplot}\NormalTok{()}
\end{Highlighting}
\end{Shaded}

\begin{center}\includegraphics[width=0.75\linewidth]{Report_files/figure-latex/page-15-1} \end{center}

\begin{Shaded}
\begin{Highlighting}[]
\FunctionTok{components}\NormalTok{(dcmptlb) }\SpecialCharTok{|\textgreater{}} \FunctionTok{autoplot}\NormalTok{()}
\end{Highlighting}
\end{Shaded}

\begin{center}\includegraphics[width=0.75\linewidth]{Report_files/figure-latex/page-15-2} \end{center}

The STL decompositions for both TFR and TLB show an overall decreasing
trend. However, there is a temporary increase beginning around 1980,
reaching a broad peak around 1991 before decreasing again. The remainder
components for both series display similar behaviour, with noticeable
dip around 1986 followed by a spike around 1988.

\newpage

\begin{Shaded}
\begin{Highlighting}[]
\NormalTok{BirthAndFertility\_MA\_TFR }\OtherTok{\textless{}{-}}\NormalTok{ BirthAndFertility }\SpecialCharTok{|\textgreater{}}
  \FunctionTok{mutate}\NormalTok{(}\StringTok{\textasciigrave{}}\AttributeTok{5{-}MA\_TFR}\StringTok{\textasciigrave{}} \OtherTok{=}\NormalTok{ slider}\SpecialCharTok{::}\FunctionTok{slide\_dbl}\NormalTok{(}\StringTok{\textasciigrave{}}\AttributeTok{Total Fertility Rate (TFR)}\StringTok{\textasciigrave{}}\NormalTok{, mean,}
                                  \AttributeTok{.before =} \DecValTok{2}\NormalTok{, }\AttributeTok{.after =} \DecValTok{2}\NormalTok{, }\AttributeTok{.complete =} \ConstantTok{TRUE}\NormalTok{))}
\NormalTok{BirthAndFertility\_MA\_TFR }\SpecialCharTok{|\textgreater{}}
  \FunctionTok{autoplot}\NormalTok{(}\StringTok{\textasciigrave{}}\AttributeTok{Total Fertility Rate (TFR)}\StringTok{\textasciigrave{}}\NormalTok{) }\SpecialCharTok{+}
  \FunctionTok{geom\_line}\NormalTok{(}\FunctionTok{aes}\NormalTok{(}\AttributeTok{y =} \StringTok{\textasciigrave{}}\AttributeTok{5{-}MA\_TFR}\StringTok{\textasciigrave{}}\NormalTok{), }\AttributeTok{colour =} \StringTok{"\#D55E00"}\NormalTok{) }\SpecialCharTok{+}
  \FunctionTok{labs}\NormalTok{(}\AttributeTok{title =} \StringTok{"Total Fertility Rate (black) and 5{-}MA estimate of trend{-}cycle (orange)"}\NormalTok{)}
\end{Highlighting}
\end{Shaded}

\begin{verbatim}
## Warning: Removed 4 rows containing missing values or values outside the scale range
## (`geom_line()`).
\end{verbatim}

\pandocbounded{\includegraphics[keepaspectratio]{Report_files/Report_files/figure-latex/page-16-1.pdf}}
The raw TFR series is already relatively smooth, so the 5-moving average
closely follows the overall trend-cycle. However, it does fail to
capture short-term fluctuations such as the dip around 1983 and spike
around 1984.

\newpage

\begin{Shaded}
\begin{Highlighting}[]
\NormalTok{BirthAndFertility\_MA\_TLB }\OtherTok{\textless{}{-}}\NormalTok{ BirthAndFertility }\SpecialCharTok{|\textgreater{}}
  \FunctionTok{mutate}\NormalTok{(}\StringTok{\textasciigrave{}}\AttributeTok{5{-}MA\_TLB}\StringTok{\textasciigrave{}} \OtherTok{=}\NormalTok{ slider}\SpecialCharTok{::}\FunctionTok{slide\_dbl}\NormalTok{(}\StringTok{\textasciigrave{}}\AttributeTok{Total Live{-}Births}\StringTok{\textasciigrave{}}\NormalTok{, mean,}
                                  \AttributeTok{.before =} \DecValTok{2}\NormalTok{, }\AttributeTok{.after =} \DecValTok{2}\NormalTok{, }\AttributeTok{.complete =} \ConstantTok{TRUE}\NormalTok{))}
\NormalTok{BirthAndFertility\_MA\_TLB }\SpecialCharTok{|\textgreater{}}
  \FunctionTok{autoplot}\NormalTok{(}\StringTok{\textasciigrave{}}\AttributeTok{Total Live{-}Births}\StringTok{\textasciigrave{}}\NormalTok{) }\SpecialCharTok{+}
  \FunctionTok{geom\_line}\NormalTok{(}\FunctionTok{aes}\NormalTok{(}\AttributeTok{y =} \StringTok{\textasciigrave{}}\AttributeTok{5{-}MA\_TLB}\StringTok{\textasciigrave{}}\NormalTok{), }\AttributeTok{colour =} \StringTok{"\#D55E00"}\NormalTok{) }\SpecialCharTok{+}
  \FunctionTok{labs}\NormalTok{(}\AttributeTok{title =} \StringTok{"Total Live{-}Births (black) and 5{-}MA estimate of trend{-}cycle (orange)"}\NormalTok{)}
\end{Highlighting}
\end{Shaded}

\begin{verbatim}
## Warning: Removed 4 rows containing missing values or values outside the scale range
## (`geom_line()`).
\end{verbatim}

\pandocbounded{\includegraphics[keepaspectratio]{Report_files/Report_files/figure-latex/page-17-1.pdf}}
The trend-cycle smooths out the short-term fluctuations in the raw TLB
series. Increasing the window size of the current 5-moving average would
result in a smoother curve.

\newpage

\begin{Shaded}
\begin{Highlighting}[]
\NormalTok{BirthAndFertility\_MA\_TLB7 }\OtherTok{\textless{}{-}}\NormalTok{ BirthAndFertility }\SpecialCharTok{|\textgreater{}}
  \FunctionTok{mutate}\NormalTok{(}\StringTok{\textasciigrave{}}\AttributeTok{7{-}MA\_TLB}\StringTok{\textasciigrave{}} \OtherTok{=}\NormalTok{ slider}\SpecialCharTok{::}\FunctionTok{slide\_dbl}\NormalTok{(}\StringTok{\textasciigrave{}}\AttributeTok{Total Live{-}Births}\StringTok{\textasciigrave{}}\NormalTok{, mean,}
                                  \AttributeTok{.before =} \DecValTok{3}\NormalTok{, }\AttributeTok{.after =} \DecValTok{3}\NormalTok{, }\AttributeTok{.complete =} \ConstantTok{TRUE}\NormalTok{))}
\NormalTok{BirthAndFertility\_MA\_TLB7 }\SpecialCharTok{|\textgreater{}}
  \FunctionTok{autoplot}\NormalTok{(}\StringTok{\textasciigrave{}}\AttributeTok{Total Live{-}Births}\StringTok{\textasciigrave{}}\NormalTok{) }\SpecialCharTok{+}
  \FunctionTok{geom\_line}\NormalTok{(}\FunctionTok{aes}\NormalTok{(}\AttributeTok{y =} \StringTok{\textasciigrave{}}\AttributeTok{7{-}MA\_TLB}\StringTok{\textasciigrave{}}\NormalTok{), }\AttributeTok{colour =} \StringTok{"\#D55E00"}\NormalTok{) }\SpecialCharTok{+}
  \FunctionTok{labs}\NormalTok{(}\AttributeTok{title =} \StringTok{"Total Live{-}Births (black) and 7{-}MA estimate of trend{-}cycle (orange)"}\NormalTok{)}
\end{Highlighting}
\end{Shaded}

\begin{verbatim}
## Warning: Removed 6 rows containing missing values or values outside the scale range
## (`geom_line()`).
\end{verbatim}

\pandocbounded{\includegraphics[keepaspectratio]{Report_files/Report_files/figure-latex/page-18-1.pdf}}
Evidently, the 7-moving average produces a smoother curve, at the cost
of losing the observations at the ends of the series.

\newpage

\section{Research Questions}\label{research-questions}

Can the ARIMA model accurately forecast the Singapore's Total Live
Births (TLB) and Total Fertility Rate (TFR) from 2013 to 2024 given the
data from 1960 to 2012, capturing the short-term fluctuations?

\section{Statistical Appendix}\label{statistical-appendix}

The following statistical methods were applied throughout the analysis
in the order presented below.

\subsection{Autoregressive Model (AR)}\label{autoregressive-model-ar}

\[
y_t = c + \phi_1 y_{t-1} + \phi_2 y_{t-2} + \cdots + \phi_p y_{t-p} + \varepsilon_t
\] where:

\(y_t\) is the value at time t

\(\phi_i\) are model parameters

\(\varepsilon_t\) is white noise

Models the current values using its previous values (lags)

\subsection{Residuals}\label{residuals}

\[
e_t = y_t - \hat{y}_t
\] \(y_t\) is the actual observed value at time t

\(\hat{y}_t\) is the predicted value at time t

The error which is the difference between the actual observed value and
predicted value at time t, used to check the chosen model accounts for
all the patterns

\subsection{Autocorrelation Function
(ACF)}\label{autocorrelation-function-acf}

\[
\rho_k = \frac{\text{Cov}(y_t, y_{t-k})}{\text{Var}(y_t)}
\] where:

\(\rho_k\) is the autocorrelation at lag \(k\)

Measures the correlation between a series \(y_t\) and lagged values
\(y_{t-k}\), which checks for trend, seasonality and whether residuals
are white noise

\subsection{Box-Cox Transformation}\label{box-cox-transformation}

\[
y^{(\lambda)} =
\begin{cases}
\frac{y^\lambda - 1}{\lambda}, & \lambda \ne 0 \\
\log(y), & \lambda = 0
\end{cases}
\]

\(\lambda\) was estimated using the Guerrero method

\subsection{STL Decomposition}\label{stl-decomposition}

\[
y_t = T_t + R_t
\] where:

\(T_t\) is the trend component,

\(R_t\) is the remainder

Decomposes time series into trend and remainder components, usually with
the seasonality component as well but the data in this analysis is
yearly

\subsection{Moving Average (MA)}\label{moving-average-ma}

\[
MA_k = \frac{1}{k} \sum_{i=-(k-1)/2}^{(k-1)/2} y_{t+i}
\] k is the window size

In this analysis when k = 5, the moving average at time t is calculated
using the 2 observations before and 2 observations after t. This process
smooths out short-term fluctuations in the series

\url{https://github.com/twalepear/TSA.git}

\end{document}
